\section{Converse Operation on a Bisupport Category}
\begin{definition} 
A compatible converse operation for a bisupport category \catcw is
an operation \converseItself such that if $f: a \morph b$ in \catcw
then $\converse{f} : b \morph a$ and satisfying:
\noindent [K.1]\quad
\[
\converse{\converse{f}}=f.
\]

\noindent [K.2]\quad
If $\sequentialdiag{a}{b}{c}{f}{g}$ in $\catc$, then
\[
\converse{(f\circ g)}
=
\converse{g}\circ\converse{f}.
\]

\noindent [K.3]\quad
For any morphism $f$,
\[
\rst{\converse{f}}=\rng{f}.
\]
\end{definition}


\begin{lemma}{Should be true.}
The category of tight zigzags is a bisupport category with compatible converse.
\end{lemma}

\begin{lemma}\commentary{\highlight{QUESTION}}
The category of tight zigzags over a bisupport category \catcw is the 
free bisupport category with converse over the bisupport category\catc. 
\end{lemma}

\subsection{Lemma that helps explain shortening}
In this lemma the morphism $s$ is a shortening.
\begin{lemma}
If \catbw is a bisupport category with compatible converse
then for all morphisms $f$, $g$ and $s$ in \catbw 
such that the following diagram commutes
\[
\downTriangleShape{x}{y_1}{y_2}
\begin{arrows}
\drawMorphism{left}{right}{a}{s}
\drawMorphism{down}{left}{a}{f}
\drawMorphism{down}{right}{b}{g}
\end{arrows}
\]
then if the determinancy condition D.2 holds of $f$ then 
\[
\converse{f} \circ g = \rng{f} \circ s. 
\]
\end{lemma}
\begin{proof}
whiteboard
\end{proof}


\subsection{Functionally-aware bisupport category}
\begin{definition}[Functionally-aware bisupport category]
A \term{function-aware bisupport category} is a bisupport
category $\catc$ together with a distinguished class of morphisms.
The distinguished morphisms form a wide subcategory $\catd\subseteq\catc$
which is closed under support and cosupport:

$$
f\in\catd
\quad\Longrightarrow\quad
\rst{f}\in\catd
\qquad\text{and}\qquad
\rng{f}\in\catd.
$$

We probably need to add R.4 for distinguished morphisms and RR.5 which should hold for all morphisms.

The distinguished morphisms represent binary relationships which are
\term{function-like}\footnote{Inclusive of partial functions.} --- those referred to in data modelling as many-one and those sometimes referred to as deterministic.
\end{definition}

\subsection{Function-aware bisupport category with converse}

\begin{definition} 
A compatible converse operation for a function-aware bisupport category \catcw is
an operation \converseItself which is compatible with the bisupport
category \catcw (see K.1, K.2 and K.3, above)
and which in addition satifies 

\noindent [K.4]\quad
For any morphism $f$ in the distinguised subset \catd,
\[
\converse{f} \circ f = \rng{f}
\]
\end{definition}

It follows from regard of the algebraic form of the definitions that 
the forgetful functor from function-aware bisupport categories with converse to the function aware bisupport categories has a left adjoint.

If we assume that the bisupport category satisfies the additional axiom RR.5 then we can describe the free function-aware bisupport category with converse as the wide subcategory of short tight zigzags. It follows that if the function-aware bisupport category satisfies condition RR.5 then the unit of this adjunction is faithful. 

\begin{aside}
To uplift shortening into the realm of zigzags over arbitrary bisupport
categories we specify shortening of $\converse{f} \circ g$ to be applicable only when $f$ is function-like.
\end{aside}

Define the function-aware category of short tight zigzags with 
converse to be the category of short tight zigzags with the
singleton zigzags $\tuple{f}$ as distinguished, for all distinguished morphisms $f$ of \catd.  

\begin{lemma}{Should be true.}
The category of function-aware short tight zigzags is a bisupport category with compatible converse.
\end{lemma}

\begin{lemma}\commentary{\highlight{THE BIG QUESTION}}
The category of function-aware short tight zigzags over a 
function-aware bisupport category \catcw is the 
free function-aware bisupport category with converse over 
the functiona-aware bisupport category\catc. 
\end{lemma}

\begin{aside}
\highlight{Where does condition R.4 get used?}
\end{aside}

\subsection{Bisupport Category with an Involutive Inner Inverse}
\begin{aside}
Is this relevant?
\end{aside}

\begin{definition}
A bisupport category with an involutive inner inverse 
is a bisupport category with compatible converse satisfying the addition 
regularity condition
\noindent [K.4]\quad
For any morphism $f$,
\[
f \circ \converse{f} \circ f =f
\]
\end{definition}

\begin{aside}
The catgeory Rel of binary relationships is a 
bisupport category with an involutive inner inverse.

It would make sense to be taking the free 
bisupport category with an involutive inner inverse on a bisupport category
as a navigational completion of a bisupport category.
However in a category of short tight zigzags the converse is \highlight{not} an inner inverse.
\end{aside}




